% SHARD-ID: RC-08-QUANTUM-IHARA-GENERAL
% SHARD-TITLE: The Ihara--Bass Formula for an Arbitrary Kraus Family
% SHARD-SUMMARY: States and proves in Lamport structure the Ihara--Bass factorisation of the quantum Hashimoto operator for an arbitrary family of superoperators with a fixed-point-free reversal, assuming no invertibility, positivity or unitarity.
% SHARD-SUMMARY: Records the three corollaries (Kraus families and trace positivity, the inverse-paired case, which factor carries the spectrum), the numerical checks, and the provenance against Watanabe--Fukumizu and Matsuura--Ohta.
% SHARD-SUMMARY: Closes with what the general formula does and does not give: the Ramanujan reading survives only in the unitary case, where it is exactly Hastings' bound.
% SHARD-KEYWORDS: Ihara-Bass, non-backtracking, quantum Hashimoto operator, Kraus family, transfer matrix, Sylvester identity, Ramanujan
% SHARD-NOTES: notes/quantum-ihara-general.md
% SHARD-SCRIPTS: scripts/qihara.py, scripts/qihara_general.py

\section{The Ihara--Bass Formula for an Arbitrary Kraus Family}
\label{sec:quantum-ihara}

This is the load-bearing section of the book. The classical Ihara--Bass formula (\cref{def:ihara-bass}) compresses a determinant on the edge
space of a regular graph to one on the vertex space, and \cref{def:quantum-hashimoto} lifts the edge operator to superoperators. The
question the founding session left open, and which TJO steered towards, is whether the compression survives when the Kraus operators are not
unitary. It does, at the price of two $\uz$-dependent operators in place of the adjacency matrix and the degree.

\emph{Status.} \Cref{thm:qihara-general} and its three corollaries are tagged \texttt{proved}. The author is
\texttt{claude:\allowbreak fable-5.1}; the reviewers are of a different model, \texttt{claude:opus}; and their reports in
\texttt{notes/reviews/} --- a round-1 algebra report, a round-1 provenance report, and a round-2 re-verification --- return
\texttt{VALID} on every verdict. Only one model family was declared for this campaign, so \texttt{proved} here means exactly that:
author $\ne$ reviewer, inside one family. The two \texttt{numerical} rows are independent evidence.

\subsection{Setting and conventions}

\begin{itemize}
  \item $\Vsp$ is a finite-dimensional complex vector space, $N = \dim\Vsp$. In the application $\Vsp = \Mn$ with the Hilbert--Schmidt inner
  product, so $N = n^2$.
  \item $\Dk = 2\mpairs$ is even, and $\rev$ is a fixed-point-free involution on $[\Dk] = \{1,\dots,\Dk\}$, the \emph{reversal}; we write
  $\bari{i} = \rev(i)$, so that $\bari{\bari{i}} = i$.
  \item $\Eop{1},\dots,\Eop{\Dk}\in\operatorname{End}(\Vsp)$ are arbitrary linear maps. No invertibility, positivity or unitarity is assumed.
  \item The edge space is $\Wsp = \Vsp\otimes\mathbb C^{\Dk}$, with $\mathbb C^{\Dk}$ spanned by $|1\rangle,\dots,|\Dk\rangle$; an element is
  $x = \sum_i x_i\otimes|i\rangle$ with $x_i\in\Vsp$ its \emph{component} at $i$.
  \item $\Hash\in\operatorname{End}(\Wsp)$ is the quantum Hashimoto operator of \cref{def:quantum-hashimoto}, $\Hash(v\otimes|i\rangle) =
  \sum_{j\ne\bari{i}}\Eop{i}v\otimes|j\rangle$, and $\zetaT(\uz) = \det\nolimits_{\Wsp}(1-\uz\Hash)^{-1}$, a rational function of
  $\uz\in\mathbb C$.
  \item $\|\cdot\|$ is the operator norm; any submultiplicative norm on $\operatorname{End}(\Vsp)$ would do. Determinants carry the space they
  are taken over, $\det\nolimits_{\Vsp}$ and $\det\nolimits_{\Wsp}$.
  \item Sylvester's identity $\det(1-XY) = \det(1-YX)$, for $X\colon\Vsp\to\Wsp$ and $Y\colon\Wsp\to\Vsp$, is used without further comment; so
  is the block determinant with identity diagonal blocks, $\det\left(\begin{smallmatrix}1 & B\\ C & 1\end{smallmatrix}\right) = \det(1-CB) =
  \det(1-BC)$ (Schur complement, then Sylvester).
\end{itemize}

\noindent\textbf{The Kraus case.} $\Vsp = \Mn$; $\Kr{1},\dots,\Kr{\mpairs}\in\Mn$ are arbitrary, $\Kr{\mpairs+k} = \Kr{k}^\dagger$, $\rev(k)
= \mpairs+k$, and $\Eop{i} = \Ad(\Kr{i})\colon\dm\mapsto \Kr{i}\dm\Kr{i}^\dagger$. With column stacking (\cref{sec:conventions}), $\Ad(A) =
\bar A\otimes A$ as an $n^2\times n^2$ matrix, so $\Tr_{\Vsp}\Ad(A) = |\Tr A|^2$ and $\Ad(A)\Ad(B) = \Ad(AB)$.

\noindent\textbf{Auxiliary maps.} Three maps do all the work: $\Rmap$ propagates along an edge, $\Smap$ starts along
every edge, $\Jrev$ is the backtracking term.
\[
  \Rmap\colon\Wsp\to\Vsp,\ \ \Rmap(v\otimes|i\rangle) = \Eop{i}v;\qquad
  \Smap\colon\Vsp\to\Wsp,\ \ \Smap v = \sum_{j=1}^{\Dk} v\otimes|j\rangle;\qquad
  \Jrev(v\otimes|i\rangle) = \Eop{i}v\otimes|\bari{i}\rangle .
\]
The two $\uz$-dependent operators on $\Vsp$ are the \emph{deformed adjacency} and the \emph{deformed degree},
\[
  \Aop \;=\; \uz\sum_{i=1}^{\Dk}\Eop{i}\,\big(1-\uz^2\Eop{\bari{i}}\Eop{i}\big)^{-1},\qquad
  \Dop \;=\; \uz^2\sum_{i=1}^{\Dk}\Eop{\bari{i}}\Eop{i}\,\big(1-\uz^2\Eop{\bari{i}}\Eop{i}\big)^{-1}.
\]
The zeta variable is $\uz$ throughout (\cref{sec:conventions}); the channel normalisation $\Sig = \sum_i\Eop{i} =
\Dk\chan$ is the one fixed there, and nothing below depends on it.

\subsection{Statement}

\begin{theorem}[Ihara--Bass for an arbitrary family of superoperators]
\label{thm:qihara-general}
\claimstatus{thm:qihara-general}{proved}
Let $\uz\in\mathbb C$ be such that $1-\uz^2\Eop{\bari{i}}\Eop{i}$ is invertible for every $i\in[\Dk]$. Then
\[
  \det\nolimits_{\Wsp}(1-\uz\Hash) \;=\;
  \prod_{\{i,\bari{i}\}}\det\nolimits_{\Vsp}\!\big(1-\uz^2\Eop{\bari{i}}\Eop{i}\big)\;\cdot\;
  \det\nolimits_{\Vsp}\!\big(1+\Dop-\Aop\big),
\]
the product being over the $\mpairs$ unordered pairs $\{i,\bari{i}\}$, which is well defined because
$\det(1-\uz^2\Eop{\bari{i}}\Eop{i}) = \det(1-\uz^2\Eop{i}\Eop{\bari{i}})$. The hypothesis holds for all
$|\uz| < (\max_i\|\Eop{\bari{i}}\Eop{i}\|)^{-1/2}$, and for all $\uz$ if every $\Eop{\bari{i}}\Eop{i} = 0$; the
identity holds pointwise at every such $\uz$, and since the left side is a polynomial of degree $\le N\Dk$ and the
right side is rational, the two sides agree as rational functions.
\end{theorem}

\begin{corollary}[Kraus families and positivity of side A]
\label{cor:qihara-kraus}
\claimstatus{cor:qihara-kraus}{proved}
For arbitrary $\Kr{1},\dots,\Kr{\mpairs}\in\Mn$ with the adjoint pairing above,
$\det\nolimits_{\Wsp}(1-\uz\Hash) = \prod_{k=1}^{\mpairs}\det(1-\uz^2\Ad(\Kr{k}^\dagger \Kr{k}))\cdot
\det(1+\Dop-\Aop)$, valid for $|\uz| < \min_k\|\Kr{k}\|^{-2}$; here $\Eop{\bari{i}}\Eop{i}$ is
$\Ad(\Kr{k}^\dagger \Kr{k})$ for $i = k\le\mpairs$ and $\Ad(\Kr{k}\Kr{k}^\dagger)$ for $i = \mpairs+k$, both positive
semidefinite. Moreover, for every $\wl\ge1$,
\[
  \Tr\nolimits_{\Wsp}\Hash^{\wl} \;=\;
  \sum_{\substack{(i_1,\dots,i_{\wl})\in[\Dk]^{\wl}\\ i_{k+1}\ne\bari{i_k}\ (k\in\mathbb Z/\wl)}}
  \big|\Tr\big(\Kr{i_{\wl}}\cdots \Kr{i_1}\big)\big|^2 \;\ge\;0,
\]
so $\zetaT(\uz) = \exp\big(\sum_{\wl}\uz^{\wl}\Tr\Hash^{\wl}/\wl\big)$ has nonnegative Taylor coefficients, and
$\zetaT(\uz) = \prod_{\pcycle}\det\nolimits_{\Vsp}(1-\uz^{\wl(\pcycle)}\Ad(\Kr{\pcycle}))^{-1}$ over primitive
cyclically non-backtracking classes $\pcycle$ of length $\wl(\pcycle)$, with $\Kr{\pcycle}$ the ordered product
along $\pcycle$.
\end{corollary}

\begin{corollary}[Inverse pairing: the unitary case]
\label{cor:qihara-unitary}
\claimstatus{cor:qihara-unitary}{proved}
If $\Eop{\bari{i}} = \Eop{i}^{-1}$ for all $i$ --- in particular for unitary Kraus operators with
$\Un{\bari{i}} = \Un{i}^\dagger = \Un{i}^{-1}$ --- then, with $\Sig = \sum_i\Eop{i}$,
\[
  \det\nolimits_{\Wsp}(1-\uz\Hash) \;=\;
  (1-\uz^2)^{N(\Dk-2)/2}\,\det\nolimits_{\Vsp}\!\big(1-\uz\,\Sig+(\Dk-1)\uz^2\big).
\]
For $\Vsp = \Mn$ and $\Sig = \Dk\chan$ this is the formula recorded as item~1 of \texttt{HANDOFF.md} and, on the
bouquet, equation~(2.11) of Matsuura--Ohta \cite{matsuura2022km}.
\end{corollary}

\begin{corollary}[Which factor carries the spectrum]
\label{cor:qihara-poles}
\claimstatus{cor:qihara-poles}{proved}
The poles of $\zetaT$ lie among the roots of $\det\nolimits_{\Vsp}(1+\Dop-\Aop) = 0$ together with the points
$\uz^2 = 1/\lambda$ for $\lambda\in\operatorname{spec}(\Eop{\bari{i}}\Eop{i})\setminus\{0\}$; the latter may cancel
against poles of the second factor, as they do in \cref{cor:qihara-unitary} (exponent $N\Dk/2$ against $-N$, a total
cancellation when $\Dk = 2$). For $\Dk = 2$ in general, $\Hash = \Eop{1}\oplus\Eop{2}$ on $\Wsp = \Vsp\oplus\Vsp$,
since the only non-backtracking successor of $i$ is $i$ itself, so $\det\nolimits_{\Wsp}(1-\uz\Hash) =
\det\nolimits_{\Vsp}(1-\uz\Eop{1})\det\nolimits_{\Vsp}(1-\uz\Eop{2})$. Conversely a pole of the second factor, which
can only sit where some $1-\uz^2\Eop{\bari{i}}\Eop{i}$ is singular, is not a pole of $\zetaT$ unless it survives
multiplication by the first.
\end{corollary}

\subsection{Proof of the theorem}

The proof is written in Lamport structure: each step is a claim followed by its own proof, and the last step assembles them.

\begin{proof}[Proof of \cref{thm:qihara-general}]
\textsc{Assume}: $\uz\in\mathbb C$ with $1-\uz^2\Eop{\bari{i}}\Eop{i}$ invertible for all $i\in[\Dk]$.
\textsc{Prove}: the displayed identity.
\begin{enumerate}
\item[$\langle1\rangle$1.] \textbf{The decomposition $\Hash = \Smap\Rmap-\Jrev$.} \emph{Proof.} On $v\otimes|i\rangle$,
$\Smap\Rmap(v\otimes|i\rangle) = \sum_j\Eop{i}v\otimes|j\rangle$ and $\Jrev(v\otimes|i\rangle) = \Eop{i}v\otimes|\bari{i}\rangle$, so
$(\Smap\Rmap-\Jrev)(v\otimes|i\rangle) = \sum_{j\ne\bari{i}}\Eop{i}v\otimes|j\rangle = \Hash(v\otimes|i\rangle)$. Both sides are linear and
these vectors span $\Wsp$. $\square$
\item[$\langle1\rangle$2.] \textbf{Block structure of $\Jrev$.} For each pair $p = \{i,\bari{i}\}$ put
$\Wsp_p = \Vsp\otimes\operatorname{span}\{|i\rangle,|\bari{i}\rangle\}$. Then $\Wsp = \bigoplus_p\Wsp_p$, $\Jrev\Wsp_p\subseteq\Wsp_p$, and
in the ordered decomposition $\Wsp_p = \Vsp\oplus\Vsp$ (component at $i$, component at $\bari{i}$), $\Jrev|_{\Wsp_p} =
\left(\begin{smallmatrix}0 & \Eop{\bari{i}}\\ \Eop{i} & 0\end{smallmatrix}\right)$. \emph{Proof.} $\Jrev$ sends the component at $i$ to the
component at $\bari{i}$ through $\Eop{i}$, and the component at $\bari{i}$ to the component at $i$ through $\Eop{\bari{i}}$, using
$\bari{\bari{i}} = i$. The direct sum is the partition of $[\Dk]$ into $\rev$-orbits, which have size two because $\rev$ is
fixed-point-free. $\square$
\item[$\langle1\rangle$3.] \textbf{$1+\uz\Jrev$ is invertible, with $\det\nolimits_{\Wsp}(1+\uz\Jrev) =
\prod_p\det\nolimits_{\Vsp}(1-\uz^2\Eop{\bari{i}}\Eop{i})$.} \emph{Proof.} By $\langle1\rangle$2, $\det\nolimits_{\Wsp}(1+\uz\Jrev) =
\prod_p\det\left(\begin{smallmatrix}1 & \uz\Eop{\bari{i}}\\ \uz\Eop{i} & 1\end{smallmatrix}\right) =
\prod_p\det\nolimits_{\Vsp}(1-\uz^2\Eop{i}\Eop{\bari{i}}) = \prod_p\det\nolimits_{\Vsp}(1-\uz^2\Eop{\bari{i}}\Eop{i})$, by the block
determinant and then Sylvester. Each factor is nonzero by \textsc{Assume}. $\square$
\item[$\langle1\rangle$4.] \textbf{Explicit inverse on a
pair.} Write $B = \uz\Eop{\bari{i}}$, $C = \uz\Eop{i}$, $P = (1-BC)^{-1} = (1-\uz^2\Eop{\bari{i}}\Eop{i})^{-1}$ and $Q = (1-CB)^{-1} =
(1-\uz^2\Eop{i}\Eop{\bari{i}})^{-1}$, both existing by \textsc{Assume} (the second is the hypothesis at the index $\bari{i}$). Then
$(1+\uz\Jrev)^{-1}|_{\Wsp_p} = \left(\begin{smallmatrix}P & -BQ\\ -CP & Q\end{smallmatrix}\right)$. \emph{Proof.} Multiply out:
$\left(\begin{smallmatrix}1&B\\C&1\end{smallmatrix}\right) \left(\begin{smallmatrix}P&-BQ\\-CP&Q\end{smallmatrix}\right) =
\left(\begin{smallmatrix}(1-BC)P & -BQ+BQ\\ CP-CP & (1-CB)Q\end{smallmatrix}\right) =
\left(\begin{smallmatrix}1&0\\0&1\end{smallmatrix}\right)$. A right inverse of a square operator on a finite-dimensional space is the
inverse. $\square$
\item[$\langle1\rangle$5.] \textbf{Sylvester step.} $\det\nolimits_{\Wsp}(1-\uz\Hash) =
\det\nolimits_{\Wsp}(1+\uz\Jrev)\cdot\det\nolimits_{\Vsp}(1-\uz\,\Rmap(1+\uz\Jrev)^{-1}\Smap)$. \emph{Proof.} By $\langle1\rangle$1,
$1-\uz\Hash = 1+\uz\Jrev-\uz\Smap\Rmap = (1+\uz\Jrev)(1-\uz(1+\uz\Jrev)^{-1}\Smap\Rmap)$, the inverse existing by $\langle1\rangle$3. Take
determinants and apply Sylvester with $X = \uz(1+\uz\Jrev)^{-1}\Smap\colon\Vsp\to\Wsp$ and $Y = \Rmap\colon\Wsp\to\Vsp$. $\square$
\item[$\langle1\rangle$6.] \textbf{Computation of $M(\uz) := \Rmap(1+\uz\Jrev)^{-1}\Smap$.} It equals
$\sum_{i=1}^{\Dk}(\Eop{i}-\uz\,\Eop{\bari{i}}\Eop{i})(1-\uz^2\Eop{\bari{i}}\Eop{i})^{-1}$. \emph{Proof.} For $v\in\Vsp$, $\Smap v$ has
component $v$ at every index. By $\langle1\rangle$4, on the pair $p = \{i,\bari{i}\}$ the vector $(1+\uz\Jrev)^{-1}\Smap v$ has component
$(P-BQ)v$ at $i$ and $(-CP+Q)v$ at $\bari{i}$. Applying $\Rmap$, which applies $\Eop{i}$ to the component at $i$, $\Eop{\bari{i}}$ to the
component at $\bari{i}$, and sums over all indices, and substituting $B,C,P,Q$, the pair $p$ contributes
\[
  \Eop{i}P-\uz\,\Eop{i}\Eop{\bari{i}}Q \;+\; \Eop{\bari{i}}Q-\uz\,\Eop{\bari{i}}\Eop{i}P
  \;=\; (\Eop{i}-\uz\Eop{\bari{i}}\Eop{i})P \;+\; (\Eop{\bari{i}}-\uz\Eop{i}\Eop{\bari{i}})Q .
\]
The first summand is the $i$-th term of the claimed sum and the second is the $\bari{i}$-th, since $Q =
(1-\uz^2\Eop{\bari{\bari{i}}}\Eop{\bari{i}})^{-1}$. Summing over the $\mpairs$ pairs gives the sum over all $\Dk$ indices. $\square$
\item[$\langle1\rangle$7.] \textbf{$1-\uz M(\uz) = 1+\Dop-\Aop$.} \emph{Proof.} Expand $\langle1\rangle$6 linearly: $\uz M(\uz) =
\uz\sum_i\Eop{i}(1-\uz^2\Eop{\bari{i}}\Eop{i})^{-1} - \uz^2\sum_i\Eop{\bari{i}}\Eop{i}(1-\uz^2\Eop{\bari{i}}\Eop{i})^{-1} = \Aop-\Dop$.
$\square$
\item[$\langle1\rangle$8.] \textbf{QED.} \emph{Proof.} Combine $\langle1\rangle$5, $\langle1\rangle$3 and $\langle1\rangle$7. Steps
$\langle1\rangle$1--$\langle1\rangle$7 hold pointwise at every $\uz$ satisfying \textsc{Assume}, with no analyticity used. The bound $|\uz|
< (\max_i\|\Eop{\bari{i}}\Eop{i}\|)^{-1/2}$ makes $\|\uz^2\Eop{\bari{i}}\Eop{i}\| < 1$, so the Neumann series supplies the inverses and
\textsc{Assume} holds there. The set of admissible $\uz$ is cofinite: its complement is the zero set of the polynomial
$\prod_i\det\nolimits_{\Vsp}(1-\uz^2\Eop{\bari{i}}\Eop{i})$, which equals $1$ at $\uz = 0$ and so is not identically zero. Two rational
functions agreeing on a cofinite set are equal. $\square$
\end{enumerate}
\end{proof}

\subsection{Proofs of the corollaries}

\begin{proof}[Proof of \cref{cor:qihara-kraus}]
$\Eop{\bari{i}}\Eop{i} = \Ad(\Kr{\bari{i}})\Ad(\Kr{i}) = \Ad(\Kr{\bari{i}}\Kr{i})$, which is
$\Ad(\Kr{k}^\dagger \Kr{k})$ or $\Ad(\Kr{k}\Kr{k}^\dagger)$; for $P\ge0$ with eigenvalues $p_a$, $\Ad(P) = \bar
P\otimes P$ has spectrum $\{\bar p_a p_b\}\subseteq[0,\infty)$, and $\|\Ad(\Kr{k}^\dagger \Kr{k})\| = \|\Kr{k}\|^4$, so
the hypothesis of \cref{thm:qihara-general} holds for $|\uz| < \min_k\|\Kr{k}\|^{-2}$; the two orderings of a pair have
equal determinant by Sylvester, leaving one factor per $k$. For the trace formula, induction on $\wl$ gives
$\Hash^{\wl}(v\otimes|i_1\rangle) = \sum\Eop{i_{\wl}}\cdots\Eop{i_1}v\otimes|i_{\wl+1}\rangle$, summed over
$(i_2,\dots,i_{\wl+1})$ with $i_{k+1}\ne\bari{i_k}$ for $k = 1,\dots,\wl$. Tracing over
$\Wsp = \Vsp\otimes\mathbb C^{\Dk}$ sets $i_{\wl+1} = i_1$, which adds the cyclic condition
$i_1\ne\bari{i_{\wl}}$, and leaves $\Tr_{\Vsp}\Ad(\Kr{i_{\wl}}\cdots \Kr{i_1}) =
|\Tr(\Kr{i_{\wl}}\cdots \Kr{i_1})|^2$. Nonnegativity of the coefficients of
$\log\zetaT = \sum_{\wl}\uz^{\wl}\Tr\Hash^{\wl}/\wl$ implies nonnegativity for $\zetaT = \exp(\cdot)$. The Euler
product is the standard regrouping of cyclic words into powers of primitive classes, using
$\sum_{k\ge1}\frac{\uz^{k\wl(\pcycle)}}{k}\Tr\Ad(\Kr{\pcycle})^k = -\log\det(1-\uz^{\wl(\pcycle)}\Ad(\Kr{\pcycle}))$,
valid for $|\uz|$ small.
\end{proof}

\begin{proof}[Proof of \cref{cor:qihara-unitary}]
With $\Eop{\bari{i}}\Eop{i} = 1$ each pair factor is $\det\nolimits_{\Vsp}(1-\uz^2) = (1-\uz^2)^N$, so the product is
$(1-\uz^2)^{N\Dk/2}$; and $\Aop = \frac{\uz}{1-\uz^2}\Sig$, $\Dop = \frac{\Dk\uz^2}{1-\uz^2}$. Hence
$1+\Dop-\Aop = \frac{1}{1-\uz^2}(1-\uz\Sig+(\Dk-1)\uz^2)$, whose determinant is
$(1-\uz^2)^{-N}\det\nolimits_{\Vsp}(1-\uz\Sig+(\Dk-1)\uz^2)$. The total exponent is $N\Dk/2-N = N(\Dk-2)/2$. The
computation divides by $1-\uz^2$ and so is carried out for $\uz^2\ne1$; both sides are polynomials in $\uz$, so the
identity extends to $\uz^2 = 1$.
\end{proof}

\begin{proof}[Proof of \cref{cor:qihara-poles}]
Immediate from the factorisation: a zero of the polynomial $\det\nolimits_{\Wsp}(1-\uz\Hash)$ is a zero of the product,
each factor is analytic away from the points named, and a polynomial has no poles, so every pole of the second factor
is cancelled by a zero of the first.
\end{proof}

\subsection{Provenance and relation to the literature}

Watanabe and Fukumizu \cite{watanabe2011zeta} prove an Ihara--Bass formula for graph zeta functions with arbitrary matrix weights, in the
shape $\det(I+\hat{\mathcal D}-\hat{\mathcal A})\prod_{[e]}\det(I-u_eu_{\bar e})$, with $\hat{\mathcal D}$ and $\hat{\mathcal A}$ built from
$(I-u_eu_{\bar e})^{-1}$ exactly as $\Dop$ and $\Aop$ are; the byte-verified quotes are in \cref{sec:prior-art}. On a one-vertex graph whose
$2\mpairs$ directed edges are $[\Dk]$ with reversal $\rev$, and with $u_e = \uz\Eop{e}$, their right-hand side becomes the right-hand side
of \cref{thm:qihara-general} after the push-through identity $\Eop{i}(1-\uz^2\Eop{\bari{i}}\Eop{i})^{-1} =
(1-\uz^2\Eop{i}\Eop{\bari{i}})^{-1}\Eop{i}$ and the relabelling $i\leftrightarrow\bari{i}$ in the degree term. Two further conventions
differ and wash out of $\det(1-\uz\,\cdot)$: they compose a prime cycle in the order opposite to our $\Eop{i_{\wl}}\cdots\Eop{i_1}$, and
their edge operator weights the target edge where $\Hash$ weights the source. Reversal of words is a bijection of cyclically
non-backtracking words because $\rev$ is an involution, and $\det(1-XY) = \det(1-YX)$.

The framing matters, and was corrected by the round-1 provenance review: their corollary is stated for loopless graphs (their graphs are
hypergraphs with two-element hyperedges), so a bouquet is not one of their graphs and their corollary does not apply to it.
\Cref{thm:qihara-general} is therefore the loop-admitting bouquet analogue of their arbitrary-weight corollary, not an instance of it; their
proof goes through a hypergraph theorem, the proof above is direct. The shape-matching itself is checked numerically only
(\texttt{notes/reviews/scratch\_wf\_mo\_conventions.py}).

Matsuura and Ohta \cite{matsuura2022km} treat invertible edge weights on a general graph, which is \cref{cor:qihara-unitary} in that
generality; on a bouquet with $X_e = \Un{e}\otimes\Un{e}^\dagger$ their equation~(2.11) is the formula of \texttt{HANDOFF.md} item~1. No
local source covers arbitrary weights on a bouquet.

\subsection{Numerical checks}

\begin{numerical}[Ihara--Bass for random unitary Kraus sets]
\label{num:qihara-unitary}
\claimstatus{num:qihara-unitary}{numerical}
For Haar-random unitary Kraus sets closed under adjoint, \cref{cor:qihara-unitary} holds at $\uz\in\{0.13,\,0.29+0.11i,\,-0.4\}$ with
relative error at most $2.1\cdot10^{-15}$ for $(n,\Dk) = (3,4)$ and at most $3.9\cdot10^{-15}$ for $(n,\Dk) = (4,6)$. The four largest
eigenvalues of $\Sig = \Dk\chan$ are $1.465,\,1.850,\,3.622,\,4.000$ for $(3,4)$ against the Ramanujan bound $2\sqrt{\qs} = 3.464$, and
$2.782,\,3.128,\,3.594,\,6.000$ for $(4,6)$ against $4.472$: the first instance violates the bound of \cref{def:ramanujan-channel} and the
second satisfies it. (\texttt{scripts/qihara.py}, \texttt{outputs/qihara.txt}.)
\end{numerical}

\begin{numerical}[The general formula in five configurations]
\label{num:qihara-general}
\claimstatus{num:qihara-general}{numerical}
\Cref{thm:qihara-general}, in both the $1-\uz M(\uz)$ and the displayed $1+\Dop-\Aop$ form, at three values of $\uz$ (one complex), with
maximal relative error per configuration as in the table; the trace formula of \cref{cor:qihara-kraus} for $\wl\le4$ to $5.5\cdot10^{-12}$;
and exact verification in rational arithmetic (sympy) for $(N,\Dk) = (1,4),\,(2,4),\,(2,6)$. (\texttt{scripts/qihara\_general.py},
\texttt{outputs/qihara\_general.txt}.)
\begin{center}
\begin{tabular}{lll}
\toprule
configuration & parameters & max.\ relative error \\
\midrule
Gaussian Kraus operators, adjoint pairing & $n=3$, $\Dk=4$ & $2.1\cdot10^{-15}$ \\
canonical-form MPS tensors and adjoints & $n=3$, $\Dk=4$ & $2.0\cdot10^{-15}$ \\
generic real superoperators, involution $(1\,2)(3\,4)$ & $n=3$, $\Dk=4$ & $6.7\cdot10^{-16}$ \\
single Kraus pair, with the direct-sum check & $n=3$, $\Dk=2$ & $1.4\cdot10^{-15}$ \\
unitary Kraus, against \cref{cor:qihara-unitary} & $n=3$, $\Dk=4$ & $1.3\cdot10^{-15}$ \\
\bottomrule
\end{tabular}
\end{center}
\end{numerical}

\subsection{What this does and does not give}

It gives the compression for \emph{every} matrix product state transfer matrix $\Sig = \sum_i\Ad(\Kr{i})$, with the physical
index doubled by the adjoints. The price of dropping unitarity is that the degree term $(\Dk-1)\uz^2$ becomes the operator $\Dop$ and the
adjacency term $\uz\Sig$ becomes the deformed channel $\Aop$, both depending on $\uz$ through $(1-\uz^2\Ad(\Kr{k}^\dagger \Kr{k}))^{-1}$.

It does not give a Ramanujan reading by itself. Without $\Eop{\bari{i}}\Eop{i} = 1$ there is no quadratic relation $\edgeev\edgeev' = \qs$
pairing the eigenvalues of $\Hash$ with those of $\Sig$. What replaces it, and whether canonical form (\cref{def:canonical-form}) buys
anything, is \cref{conj:kraus-ramanujan}.

Side A is unconditionally positive: $\Tr\Hash^{\wl} = \sum_{\pcycle}|\Tr \Kr{\pcycle}|^2\ge0$ for any Kraus family, so the primes of this
zeta are the primitive cyclically non-backtracking words in the $\Kr{k}$ and $\Kr{k}^\dagger$, with real nonnegative multiplicities.

\begin{observation}[RH for the quantum Ihara zeta is Hastings' bound]
\label{obs:rh-iff-ramanujan-channel}
\claimstatus{obs:rh-iff-ramanujan-channel}{sketched}
For unitary Kraus operators closed under adjoint, \cref{cor:qihara-unitary} attaches to each eigenvalue $\lambda$ of $\Sig = \Dk\chan$ the
quadratic $1-\lambda\uz+\qs\uz^2 = (1-\edgeev\uz)(1-\edgeev'\uz)$ with $\edgeev\edgeev' = \qs = \Dk-1$ and $\edgeev+\edgeev' = \lambda$; the
eigenvalues $\lambda$ are real because $\Sig$ is Hilbert--Schmidt self-adjoint for an adjoint-closed set. Both roots then have $|\edgeev| =
\sqrt{\qs}$ exactly when $|\lambda|\le 2\sqrt{\qs}$. Hence every nontrivial pole of $\zetaq$ lies on the circle $|\uz| = \qs^{-1/2}$ if and
only if $\Dk\lamtwo\le 2\sqrt{\Dk-1}$, which is Hastings' bound $\lamH$ \cite{hastings2007random} in the adjacency units of
\cref{sec:conventions}: the Riemann hypothesis for this zeta is exactly the statement that $\chan$ is a Ramanujan quantum expander
(\cref{def:ramanujan-channel}). The note records this as an observation, not found in the literature in this form; it is tagged
\texttt{sketched} because it has had no reviewer.
\end{observation}
